% ============================================================
%  CHAPITRE 1 : Comprendre le Problème avant de le Résoudre
% ============================================================
\chapter{Comprendre le Problème avant de le Résoudre}
\label{chap:contexte}

\begin{infobox}[Résumé du chapitre]
Ce chapitre pose les fondations du projet~: définir ce qu'est le MLOps et
pourquoi il est indispensable en santé, caractériser la problématique médicale,
décrire le dataset CMS~DE-SynPUF, présenter la stack technologique retenue et
illustrer l'architecture globale du pipeline.
\end{infobox}

% ============================================================
\section{Le MLOps~: industrialiser l'IA pour qu'elle tienne ses promesses}
% ============================================================

\subsection{De la recherche à la production~: un fossé difficile à franchir}

Les prévisions de la littérature académique et des rapports industriels
convergent~: \textbf{plus de 85\,\% des projets de machine learning ne parviennent
jamais en production}~\cite{sculley2015hidden, zaharia2018accelerating}. Cette
statistique, souvent citée, recouvre une réalité concrète~: un modèle entraîné
dans un notebook peut obtenir d'excellentes métriques en validation, et pourtant
échouer à produire de la valeur lorsqu'il est confronté aux flux de données réels,
aux contraintes d'infrastructure ou à l'évolution du comportement des données dans
le temps.

Le terme \textbf{MLOps} --- contraction de \textit{Machine Learning} et
\textit{Operations} --- désigne l'ensemble des pratiques, outils et processus visant
à réduire ce fossé. Il s'inspire du mouvement DevOps qui a transformé le
développement logiciel, en y ajoutant les contraintes spécifiques aux systèmes
d'apprentissage automatique~: gestion des données, versionnement des modèles,
reproductibilité des expériences et surveillance de la dérive.

\subsection{Les huit principes du MLOps}

\begin{table}[htbp]
\centering
\caption{Les huit principes fondamentaux du MLOps}
\label{tab:principes_mlops}
\begin{tabularx}{\textwidth}{@{}lX@{}}
\toprule
\textbf{Principe} & \textbf{Description} \\
\midrule
Reproductibilité & Toute expérience doit être reproductible à l'identique \\
Versionnement    & Données, code et modèles versionnés conjointement \\
Automatisation   & Pipeline de bout en bout sans intervention manuelle \\
Surveillance     & Monitoring continu des performances et de la dérive \\
Traçabilité      & Chaque prédiction est tracée avec ses paramètres \\
Scalabilité      & Architecture capable de monter en charge \\
Équité           & Absence de biais systématique par groupe protégé \\
Gouvernance      & Auditabilité et conformité réglementaire \\
\bottomrule
\end{tabularx}
\end{table}

\subsection{Niveaux de maturité MLOps}

Google~\cite{sculley2015hidden} définit trois niveaux de maturité~:
\begin{itemize}
  \item \textbf{Niveau~0} — entraînement manuel, déploiement ponctuel, pas de
        surveillance ;
  \item \textbf{Niveau~1} — pipeline ML automatisé avec réentraînement déclenché
        par les données ;
  \item \textbf{Niveau~2} — pipeline CI/CD complet, orchestration, tests
        automatisés sur le modèle.
\end{itemize}

\noindent Ce projet vise le \textbf{niveau~2}~: toute modification du code déclenche
une chaîne d'intégration continue (GitHub Actions), le monitoring surveille
automatiquement la dérive et déclenche le rollback si nécessaire.

% ============================================================
\section{Problématique médicale~: prévenir plutôt que guérir}
% ============================================================

\subsection{Le fardeau des hospitalisations évitables}

Les hospitalisations évitables constituent un problème de santé publique majeur.
Elles surviennent principalement chez les patients \textbf{polypathologiques} ---
ceux qui souffrent de plusieurs maladies chroniques simultanément. Ces patients
cumulent des consultations fréquentes, des ordonnances multiples et un risque
significativement élevé d'être hospitalisés en urgence.

Identifier ces patients \textit{avant} que leur état ne se dégrade permet
d'intervenir précocement~: ajustement thérapeutique, consultation de coordination,
programme d'éducation thérapeutique. Ces interventions coûtent en moyenne
\textbf{15~fois moins} qu'une hospitalisation d'urgence.

\subsection{L'index de Charlson~: quantifier la charge de morbidité}

L'index de Charlson~\cite{charlson1987new} est un score clinique validé qui
prédit la mortalité à un an en fonction des comorbidités. Dans notre projet, il
est calculé comme suit~:

\begin{equation}
\label{eq:charlson}
\text{CHARLSON} = 2 \cdot \text{REIN} + 2 \cdot \text{CANCER} +
2 \cdot \text{AVC} + \sum_{k \in \mathcal{C}_1} \text{SP}_k
\end{equation}

\noindent où $\mathcal{C}_1 = \{\text{CHF}, \text{DIABÈTE}, \text{BPCO},
\text{ALZHEIMER}, \text{DÉPRESSION}, \text{CORONAROPATHIE},
\text{ARTHRITE}\}$ sont les comorbidités de poids~1.

\noindent La valeur moyenne dans notre dataset est de $\overline{\text{CHARLSON}}
= 2.34$, avec un maximum observé de~11.

\subsection{Le score PSI~: surveiller la dérive des données}

Le \textbf{Population Stability Index (PSI)} mesure la dérive entre une
distribution de référence et une distribution courante~:

\begin{equation}
\label{eq:psi}
\text{PSI} = \sum_{i=1}^{n} \left( p_i - q_i \right)
             \ln\!\left(\frac{p_i}{q_i}\right)
\end{equation}

\noindent Par convention~: PSI~$< 0.10$ = stable, $0.10$--$0.20$ = attention,
$> 0.20$ = réentraînement requis.

% ============================================================
\section{Le dataset CMS DE-SynPUF~: une mine de données réelles}
% ============================================================

\subsection{Description générale}

Le \textit{CMS Data Entrepreneurs Synthetic Public Use File} (DE-SynPUF) est
publié par les Centers for Medicare \& Medicaid Services (CMS) américains. Il
constitue une base de données \textbf{synthétique mais réaliste}, construite
pour préserver les distributions statistiques des données réelles tout en
garantissant la confidentialité des patients.

\begin{table}[htbp]
\centering
\caption{Caractéristiques générales du dataset CMS DE-SynPUF}
\label{tab:dataset}
\begin{tabular}{@{}ll@{}}
\toprule
\textbf{Caractéristique} & \textbf{Valeur} \\
\midrule
Volume total         & 60~Go, 20~samples \\
Bénéficiaires totaux & 2,3~millions \\
Sample utilisé       & Sample~1 --- 30\,000 patients \\
Période temporelle   & 2008 à 2010 \\
Pays                 & États-Unis (Medicare) \\
\bottomrule
\end{tabular}
\end{table}

\subsection{Structure des fichiers}

Le dataset est structuré en cinq types de fichiers complémentaires~:

\begin{table}[htbp]
\centering
\caption{Volume des fichiers claims pour 30\,000 patients}
\label{tab:claims}
\begin{tabular}{@{}lrl@{}}
\toprule
\textbf{Fichier} & \textbf{Lignes} & \textbf{Contenu} \\
\midrule
Bénéficiaires  &   88\,585 & Données démographiques et comorbidités \\
Inpatient      &   17\,325 & Hospitalisations \\
Outpatient     &  205\,642 & Consultations ambulatoires \\
Carrier (A+B)  & 1\,233\,861 & Actes médicaux (médecins) \\
Prescription   & 1\,436\,994 & Médicaments délivrés \\
\midrule
\textbf{Total} & \textbf{\textasciitilde{}2,9~M} & \\
\bottomrule
\end{tabular}
\end{table}

\subsection{Distribution démographique}

Sur les 30\,000 patients du Sample~1~:
\begin{itemize}
  \item \textbf{Genre~:} 55.3\,\% de femmes, 44.7\,\% d'hommes ;
  \item \textbf{Groupes raciaux~:} Blanc (majoritaire), Noir, Hispanique,
        Asiatique, Natif Américain, Autre ;
  \item \textbf{490 décès} enregistrés sur la période ;
  \item Dates corrigées~: du 2007-12-10 au 2010-12-30.
\end{itemize}

% ============================================================
\section{Stack technologique~: des choix architecturaux justifiés}
% ============================================================

\subsection{Vue d'ensemble}

La stack retenue équilibre \textbf{maturité industrielle, performance et
maintenabilité}. Chaque composant répond à un besoin précis et a été préféré
à ses alternatives pour des raisons documentées.

\begin{table}[htbp]
\centering
\caption{Stack technologique du projet}
\label{tab:stack}
\begin{tabular}{@{}lll@{}}
\toprule
\textbf{Couche} & \textbf{Outil} & \textbf{Rôle} \\
\midrule
Modélisation      & XGBoost, LightGBM, Scikit-learn & Entraînement \\
Optimisation      & Optuna                          & Recherche hyperparamètres \\
Traçabilité       & MLflow                          & Expériences et registre \\
Versionnement     & DVC + Git                       & Données et modèles \\
Validation        & Great Expectations              & Qualité des données \\
Orchestration     & Prefect                         & Pipelines de données \\
Déploiement       & FastAPI + Uvicorn               & API REST \\
Conteneurisation  & Docker + Docker Compose         & Isolation \\
CI/CD             & GitHub Actions                  & Intégration continue \\
Monitoring        & Evidently AI + Grafana          & Dérive et alertes \\
\bottomrule
\end{tabular}
\end{table}

\subsection{Pourquoi pas Kafka, Spark, Redis ou Kubernetes~?}

\begin{alertbox}[Choix architecturaux délibérés]
Ces technologies ne sont pas absentes par limitation technique, mais parce
qu'elles seraient \textbf{surdimensionnées} pour ce projet~:
\begin{itemize}
  \item \textbf{Apache Kafka~:} le dataset est historique et statique --- un
        pipeline batch suffit ;
  \item \textbf{Apache Spark~:} 30K patients et 2.9M lignes s'exécutent en
        quelques secondes avec Pandas~;
  \item \textbf{Redis (Feature Store)~:} l'API répond en $<$100ms sans cache
        distribué~;
  \item \textbf{Kubernetes~:} Docker Compose couvre les besoins de déploiement ;
        Kubernetes interviendrait à l'échelle des 2.3M patients.
\end{itemize}
\end{alertbox}

% ============================================================
\section{Architecture globale du pipeline MLOps}
% ============================================================

La figure~\ref{fig:pipeline_mlops} illustre l'architecture complète du pipeline,
depuis l'ingestion des données jusqu'au monitoring en production.

\begin{figure}[htbp]
\centering
% ------------------------------------------------------------------
%  FIGURE 1.1 — Pipeline MLOps (TikZ)
% ------------------------------------------------------------------
\begin{tikzpicture}[
  font=\small,
  >=Stealth,
  box/.style={
    rectangle, rounded corners=4pt, minimum width=3cm, minimum height=1cm,
    text centered, text width=2.8cm, draw, thick
  },
  databox/.style={box, fill=lightblue,   draw=primary},
  mlbox/.style  ={box, fill=lightgreen,  draw=success},
  deploybox/.style={box, fill=lightyellow, draw=accent},
  monbox/.style ={box, fill=danger!10,   draw=danger},
  arrow/.style  ={->, thick, draw=gray80},
  label/.style  ={font=\footnotesize\bfseries, text=white, rounded corners=2pt}
]

% Row 1 — Données
\node[databox] (raw)    at (0,0)    {Données brutes\\CMS DE-SynPUF};
\node[databox] (dvc)    at (3.5,0)  {Versionnement\\DVC + Git};
\node[databox] (ge)     at (7,0)    {Validation\\Great Expectations};
\node[databox] (feat)   at (10.5,0) {Feature\\Engineering\\(29 variables)};

% Row 2 — Modèles
\node[mlbox] (train) at (0,-3)   {Entraînement\\XGBoost / LGB\\Logistique};
\node[mlbox] (optuna) at (3.5,-3) {Optimisation\\Optuna\\(50 essais)};
\node[mlbox] (mlflow) at (7,-3)  {Traçabilité\\MLflow\\Registre};
\node[mlbox] (eval)  at (10.5,-3) {Évaluation\\AUC, Recall\\SHAP, Fairness};

% Row 3 — Production
\node[deploybox] (api)    at (0,-6)   {API FastAPI\\Docker};
\node[deploybox] (cicd)   at (3.5,-6) {CI/CD\\GitHub Actions};
\node[monbox]    (evid)   at (7,-6)   {Monitoring\\Evidently AI\\PSI};
\node[monbox]    (grafana) at (10.5,-6){Dashboard\\Grafana\\+ Rollback};

% Row 1 arrows
\draw[arrow] (raw)   -- (dvc);
\draw[arrow] (dvc)   -- (ge);
\draw[arrow] (ge)    -- (feat);

% Row 1 → Row 2
\draw[arrow] (feat)  -- (10.5,-1.5) -- (0,-1.5) -- (train);

% Row 2 arrows
\draw[arrow] (train) -- (optuna);
\draw[arrow] (optuna) -- (mlflow);
\draw[arrow] (mlflow) -- (eval);

% Row 2 → Row 3
\draw[arrow] (eval)  -- (10.5,-4.5) -- (0,-4.5) -- (api);

% Row 3 arrows
\draw[arrow] (api)    -- (cicd);
\draw[arrow] (cicd)   -- (evid);
\draw[arrow] (evid)   -- (grafana);

% Feedback loop
\draw[arrow, dashed, draw=danger, thick]
  (grafana) -- ++(1.5,0) -- ++(0,7.5) node[midway, right, text=danger,
  font=\footnotesize\bfseries] {Rollback / Réentraînement}
  -- ++(-14.5,0) -- (raw);

\end{tikzpicture}
\caption{Architecture globale du pipeline MLOps --- de l'ingestion au monitoring}
\label{fig:pipeline_mlops}
\end{figure}

\subsection{Description des flux}

Le pipeline se déroule en quatre étapes séquentielles~:

\begin{enumerate}
  \item \textbf{Couche données~:} les fichiers bruts CMS sont versionnés via DVC,
        validés par Great Expectations, puis transformés en 29 features exploitables.
  \item \textbf{Couche modèles~:} les trois modèles sont entraînés avec un split
        temporel strict. Optuna optimise les hyperparamètres sur 50 essais. MLflow
        trace chaque expérience et conserve le registre des modèles.
  \item \textbf{Couche déploiement~:} le modèle champion est empaqueté dans un
        conteneur Docker et exposé via FastAPI. GitHub Actions assure que tout
        commit déclenche les tests et redéploie automatiquement si les tests passent.
  \item \textbf{Couche monitoring~:} Evidently calcule le PSI quotidiennement.
        Si PSI~$>~0.20$ ou AUC~$<~0.75$, Grafana émet une alerte et le système
        bascule automatiquement vers la version précédente du modèle (\textit{rollback}).
\end{enumerate}

\section*{Synthèse du chapitre}

Ce premier chapitre a posé le cadre conceptuel, médical et technique du projet.
Le MLOps répond à un besoin réel~: transformer un modèle de recherche en un service
fiable, traçable et équitable. Le dataset CMS~DE-SynPUF, avec ses 2.9~millions
de lignes pour 30\,000 patients, fournit une matière riche et réaliste. La stack
technologique choisie priorise la maturité et la proportionnalité aux besoins.
Le chapitre suivant détaille comment cette donnée brute a été transformée en
données prêtes à l'entraînement.
